\documentclass[10pt,a4paper,ragged2e,withhyper]{custom-altacv}

% Specify the CV version according to format: v.YY.MM.DD-L.S
% where <L> is language (RU, EN, etc.) and <S> is a special mark like 'demo'.
\input{preamble.tex}


\begin{document}

\name{Казуро Денис}
\tagline{Project Manager}

\personalinfo{%
  \email{denis.kazuro@yandex.com}
%  \phone{}
  \telegram{creppuscule}
%  \printinfo{\faLinkedin}{Фартыгин Артём}[https://linkedin.com/in/artyom-fartygin-61890024a]
 \github{Levatei}
  \location{Россия, Москва}
}

\makecvheader

%% Depending on your tastes, you may want to make fonts of itemize environments slightly smaller
% \AtBeginEnvironment{itemize}{\small}

%% Set the left/right column width ratio to 6:4.
\columnratio{0.6}

% Start a 2-column paracol. Both the left and right columns will automatically break across pages
% if things get too long.
\begin{paracol}{2}

\cvsection{Опыт Работы}

\cvjob{Project Manager }{Web3 Startup / Product Studio }{Январь 2023 — Декабрь 2025 (3 года)}
\begin{itemize}
	\item Руководил полным циклом реализации проекта: от формирования идеи и сбора требований до запуска рабочего MVP.
	\item Внедрил с нуля Scrum + Kanban, повысив предсказуемость delivery \textbf{(85-90\% выполнения спринтов)}, увеличив стабильность релизов. 
	\item Формировал и управлял двумя кросс-функциональными удалёнными \textbf{командами до 10 человек}.
    \item Внедрил системный онбординг-пакет, \textbf{сократив time-to-productivity новых сотрудников на 26\%}.
    \item Вёл бэклог проекта: формирование спринтов, приоритизация,  контроль исполнения с полным контролем ключевых метрик (\textbf{Velocity, Burndown, Cycle Time}). 
    \item \textbf{Снизил scope creep на 25-30\%} и количество пересогласований на 30\% , выстроив процесс управления требованиями с фиксацией решений после каждого звонка со стейкхолдерами.
    \item Формировал технические задания для дизайн команды, согласовывал дизайн-проекты и контролировал их реализацию.
    \item Вывел 2 ключевых продукта на первичную монетизацию за 6 месяцев с момента старта;\textbf{ ROI} одного проекта составил \textbf{1}\textbf{0X к бюджету разработки} \textbf{(40k\$). }
\end{itemize}

\divider

\cvjob{Project Manager }{Huntica }{Июль 2021 — Ноябрь 2022 (1 год и 5 месяцев) }{}{}

\begin{itemize}
	\item \textbf{Внедрил} с нуля \textbf{Agile-церемонии} (Daily, Review, Retro), повысив прозрачность процессов для фаундера и PO.
    \item \textbf{Руководил кросс-функциональной командой}: дизайнеры, front-end и back-end разработчики.
    \item Собирал, анализировал и документировал \textbf{требования пользователей и стейкхолдеров }для планирования спринтов. 
    \item Внедрил Toggl для учёта рабочего времени, что позволило оптимизировать систему оплаты труда. 
    \item \textbf{Оптимизировал Jira}, сократив время на обработку задач \textbf{на 20\%} и уменьшив количество дубликатов на 90\%.
\end{itemize}

\divider

\cvjob{Project Manager }{МВА Гарант-групп }{Апрель 2019 — Сентябрь 2019 (6 месяцев)  }{}{}
    \begin{itemize}
	\item Внедрил CRM-систему, обеспечив запуск в срок и в соответствии с требованиями. 
    \item Формировал и ставил задачи подрядчикам (web-разработка, обслуживание веб-инфраструктуры), контролировал сроки и качество выполнения.
	\item \textbf{Снизил потерю входящих заявок на 95\%} в сравнении с отчётным периодом. 
    \item Обеспечивал прозрачную коммуникацию между стейкхолдерами и техническими исполнителями. 
\end{itemize}




%%%%% Temporarily commented out
\iffalse
\cvsection{Сертификаты и Дипломы}

\cvachievement{\faGraduationCap}{Участие в двух \textbf{международных} олимпиадах по астрономии}{IOAA и IAO}
\fi
%%%%%



%% Switch to the right column. This will now automatically move to the second
%% page if the content is too long.
\switchcolumn

\cvsection{Образование}

\cvevent{\faUniversity Среднее образование}{\mbox{Минский колледж предпринимательства}}{Сентябрь, 2016 -- Июль, 2019}{}
Выпускник по направлению коммерческая деятельность.



\cvsection{Навыки}

\vspace{2ex}

\begin{itemize}
	\item \cvtag{Agile}  \cvtag{Scrum}\cvtag{Kanban} 
	\item \cvtag{Jira} \cvtag{Confluence}  \cvtag{Asana}\cvtag{Figma} 
    \item \cvtag{Miro} \cvtag{MS Excel} \cvtag{PMBOK}
	\item \cvtag{Git}  \cvtag{API}\cvtag{Postman} \cvtag{Devtools}
    \item \cvtag{Управление бэклогом} 
    \item \cvtag{Управление стейкхолдерами}

\end{itemize}

\divider

\textcolor{emphasis}{Имел опыт работы}
\vspace{2ex}

\begin{itemize}
	\item \cvtag{Python} \cvtag{HTML/CSS}
	\item  \cvtag{PostgreSQL}\cvtag{SQL Lite}
	\item  \cvtag{Claude Code} \cvtag{Cursor} \cvtag{OpenClaw}


\end{itemize}



\cvsection{Языки}

\begin{itemize}
	\item Английский --- \cvtag{Intermediate}
    \item Русский    --- \cvtag{Носитель}
\end{itemize}

\cvsection{О себе}

Project Manager  уровня Middle с опытом запуска продуктов в Web3, HR Tech и SaaS.
 
\item  Реализовал больше 10 проектов от идеи до монетизации, ориентируясь на Agile-метрики.
\item  Умею управлять паралельно несколькими распределенными командами.
\item   Благодаря техническому бэкграунду в QA, понимаю разработку на уровне API и интеграций и говорю с командой на одном языке. 
	
\end{paracol}

\end{document}
